\paragraph{Matching the image adaptation horizon does not recover robustness}
\label{sec:horizon500_outcome}
The final prespecified revision increases both clean-teacher adaptation and the student's simulated adaptation from 100 to 500 inner updates. The fixed-reference and adapted-reference models retain the same 100 outer updates and prediction-retention coefficient 32. All three supplied models pass the fresh initial screen, with 32/32 detections and preservation passes. This allows a matched attack comparison; it does not itself establish robustness.

The blind attacker uses clean data outside the defense support, with checkpoints at 25, 100 and 500 updates and prediction-anchor weights 0, 1, 4 and 16. Quality-only selection on 32 validation images chooses the 500-update, zero-anchor checkpoint for each supplied model. Independent owner evaluation then uses 64 final generation seeds and fitting-holdout classes, with the previously frozen detector and preservation thresholds. Validation and final class sets are disjoint; their identities were used in earlier development studies, so this is not new-class confirmation. No owner scores enter checkpoint selection.

\begin{table}[ht]\centering\small
\begin{tabular}{lrrrrr}\toprule
Starting model & Updates & Anchor & Preserved & Removed & Joint removed\\\midrule
Original & 500 & 0.0 & 64/64 & 20/64 & 20/64\\
Fixed & 500 & 0.0 & 64/64 & 21/64 & 21/64\\
Adapted & 500 & 0.0 & 64/64 & 20/64 & 20/64\\
\bottomrule\end{tabular}
\caption{Independent open-weight clean-training attacks against DistSeal UViT-H/DC-AE at 512 pixels. Original is the released marked generator; Fixed and Adapted are the two prediction-anchored meta-distilled variants. Updates counts clean attacker updates, and Anchor is its prediction-retention coefficient. Preserved counts returned outputs passing LPIPS $\leq0.2$, SSIM $\geq0.6$, image CLIP $\geq0.9$ and class-alignment retention $\geq0.95$. Removed counts returned images losing owner detection, before requiring preservation. Joint counts both outcomes. All fractions use 64 reserved final inputs per starting model, initially all detected at strict bit accuracy $>0.6875$. Checkpoint selection uses 32 separate validation inputs without owner feedback. This experiment returns the marked source when the raw adapted image fails preservation. The reported counts concern returned outputs; the raw adapted generator can fail the pair rule. Updates and anchor coefficients are selected separately for each starting model.}\end{table}


The quality-fallback client achieves 20/64 joint removals for the released model, 21/64 for the fixed-reference model and 20/64 for the adapted-reference model. All 64 returned images pass preservation in each case. However, this client retains the supplied generator and returns its source image whenever the raw attack fails the gates. Before fallback, raw preservation passes are only 21/64, 23/64 and 21/64, and raw removals are 58/64, 56/64 and 57/64. Consequently, its preservation rate is a property of the two-path client, not evidence that the adapted checkpoint alone preserves every source image. The clean null has 0/64 detections, which remains too small for precise low-FPR validation.

Neither training variant reduces quality-preserving removal relative to the released model. The longer adaptation horizon therefore does not rescue this meta-distillation hypothesis. Following the rejection criterion frozen before training completion, this direction stops here and is not expanded to additional watermark families. A different image defense would require a distinct mechanism, a source-grounded novelty assessment and a new experiment. These negative results do not support a generally robust image watermark.

The counts above retain the original frozen quality rule. The \hyperref[sec:quality_recalibration]{post-hoc reanalysis} applies a revised gate to the same saved pairs.
